\documentclass[10pt]{article}
\usepackage[a4paper,margin=0.95in]{geometry}
\usepackage{amsmath,amssymb}
\usepackage{booktabs}
\usepackage{graphicx}
\PassOptionsToPackage{hyphens}{url}
\usepackage[hidelinks]{hyperref}
\usepackage{url}
\usepackage{microtype}
\newcommand{\eps}{\varepsilon}

\title{Load, Granularity and Reader Effects: Exploratory Ablations of\\ $\eps$-Sufficient Context for Language-Model Code Editing in Go (Phases 1 to 3)}
\author{Jude Castiel Johansen\\ Shockware Limited, independent researcher\\ \texttt{castiel@shockware.co.uk} \quad ORCID 0009-0002-5644-3073}
\date{18 October 2026 \quad\textit{Exploratory report, version 1.0}}

\begin{document}
\maketitle

\begin{abstract}
This report describes three exploratory ablation phases run before the confirmatory study of a registered report on $\eps$-sufficient context for language-model code editing in Go. None of the results is confirmatory. Some aggregates were seen before repairs, the instance sets were chosen by screens on the evaluated models, and eleven instances from two repositories are a small basis.

Phase 1 removed dependency interfaces one group at a time on 23 pilot instances (Haiku 5.5, $k=20$). The count of apparently necessary groups (33 of 135) lay inside a null reference.

Phase 2 ran four models on 10 screened instances. On Haiku 5.5, 7/10 instances had no arm reaching 0.1. No average difference was resolved between the write-set bodies (W), W plus dependency interfaces (A1) and the predictor-selected arms.

Leave-one-out at $k=100$ gave counts inside their null ranges (on two of the three instances the test could not reject by construction). On cli 2bcdfbb1ac02, however, removing one package moved the pass rate from 0.27 to 0.00.

Phase 3 tested declaration-level arms and padding intended to match A1's prompt tokens on 11 instances with headroom. On Haiku 5.5, a reachable-declaration slice of W met the non-inferiority margin against W only under a post-hoc sample-set choice (post-screen samples, not preregistered; one-sided lower bound $-0.043$, median token ratio 0.47). It did not meet the margin with screen samples included ($-0.051$), nor with the lower end of the two-sided interval ($-0.053$). On Sonnet 5.5 its point estimate was below W's ($-0.045$ $[-0.082, -0.009]$); at $k=10$ this is an indication, not a resolved effect. Haiku 5.5 pass rates were lower with padding (Wpad$-$W $-0.138$ $[-0.281, -0.018]$); token matching fell short on two of the three caddy instances, where the same-module padding pool was exhausted.
\end{abstract}

\section{Purpose and relation to the prior publications}

This report belongs to a series of public records on $\eps$-sufficient context for language-model code editing. The definitions note \cite{definitions} defines the concept. The pilot \cite{pilot} ran 24 instances from \texttt{cli/cli}. One of these had an empty task text and is invalid (erratum of 2026-10-12), which leaves 23 valid pilot instances.

The Stage 1 protocol \cite{stage1} operationalises $\eps$-sufficiency as a paired, single-shot, multi-model test. A compiler-computed interface context is compared with bodies and with the transitive slice. The outcome is $\Delta$-pass, the pass rate of the fail-to-pass tests after the model's patch is applied.

The three phases reported here were run between 2026-10-12 and 2026-10-14. They took place after the pilot and before any confirmatory call. Their purpose was to decide which arms and screens the confirmatory study should use. They are exploratory throughout, and we report them regardless of outcome.

The preregistration states the hypotheses relevant here as follows.

\begin{itemize}
\item \textbf{H2 (primary).} A1 is non-inferior to A1b (dependency bodies) and to A2 (transitive slice) at margin $\eps$ on $\Delta$-pass, tested by intersection--union.
\item \textbf{H1 (revised after the pilot).} At each preregistered load level and padding type, the protocol tests (a) one-sided paired superiority of A1 over A1 padded to that level (degradation) and (b) TOST equivalence at $\pm 0.05$ (flatness). Each level is classed as degradation, flat or inconclusive, and the full curve is reported with intervals.
\item \textbf{H5 (added 2026-10-12, before any ablation call).} H5 concerns the necessity and static predictability of dependency interfaces within A1. It is confirmatory on ``the Phase 3 instances'' and exploratory on the 24 pilot instances.
  \begin{itemize}
  \item For each direct dependency $d$ of W, arm A1$-d$ removes $d$'s interface from A1. When there are many dependencies, they are grouped into contiguous groups.
  \item Arm A1g keeps only the dependencies referenced by the gold patch (an oracle predictor). Arm A1t keeps only those referenced by the task text (an inference-time predictor).
  \item H5a: A1g is non-inferior to A1 at $\eps$. H5b: A1t is non-inferior to A1 at $\eps$.
  \item Descriptive quantities are fixed at $k=20$: the fraction of necessary dependencies, the precision and recall of each predictor, a slack ratio, and the count of removals that raise the pass rate by more than $\eps$.
  \end{itemize}
\end{itemize}

None of H1, H2 or H5 was tested confirmatorily in these phases. H1 and H2 belong to the confirmatory run. H5 was tested exploratorily in Phases 1 and 2. Section~\ref{sec:lim} discusses what ``the Phase 3 instances'' in the H5 paragraph refers to.

The work relates to three strands of published research. The first is the notion of sufficient context in retrieval-augmented generation \cite{joren2025}. The second is the observation, stated in its title, that context length alone can hurt performance despite perfect retrieval \cite{du2025}. The third is evidence on how position within long contexts affects use of information \cite{liu2024}. We make no claim of novelty beyond the specific arms, instances and contrasts reported below.

\section{Design}

\subsection{Instances and screens}

All instances are fail-to-pass commits from single-module Go repositories. They were mined from commits since 2026-06-01; the mining windows covered 159 first-parent commits for caddy, 240 for hugo and 13 for zap (none valid). The instance directories number 14 for \texttt{caddyserver/caddy} and 30 for \texttt{cli/cli}; the latter include the 24 pilot instances (one invalid) and earlier template/test instances, so the count of valid \texttt{cli/cli} instances is lower. There is one valid instance each for \texttt{gohugoio/hugo}, \texttt{restic/restic} and \texttt{golangci/golangci-lint}. Every examined commit and its rejection reason is listed in \texttt{harness/instances/rejected.jsonl}.

\paragraph{Phase 1 (2026-10-12).} Phase 1 used the 23 valid pilot instances, all from \texttt{cli/cli}, with no screen.

\paragraph{Phase 2 (2026-10-12/13).}
\begin{itemize}
\item \emph{Screen.} We screened 35 instances (22 cli, 10 caddy, 1 hugo, 1 restic, 1 golangci-lint) with arm W at $k=5$ on Haiku 5.5. Under the amendment of 2026-10-12, an instance entered only if W was inexpressible (a gold identifier absent from W) and W passed below 0.8. The results file records 162 screen samples, fewer than 35 instances at $k=5$ would give. The shortfall is not explained in the results file.
\item \emph{Selection.} Sixteen instances had W inexpressible, and 10 were selected (6 cli, 4 caddy).
\item \emph{Missing instance.} Instance cli bd81afda9ab6 is absent because its arms failed to build during a disk-full interruption.
\end{itemize}

\paragraph{Phase 3 (2026-10-13/14).}
\begin{itemize}
\item \emph{Screen.} We screened 35 instances with W and A1 at $k=20$ on Haiku 5.5. An instance entered if either rate lay in $[0.1, 0.9]$.
\item \emph{Selection.} Twelve instances were selected, including \texttt{gohugoio/hugo} c9f6e4d4abc3. That instance was dropped on 2026-10-14 because it produced a 182 GB build cache. The preregistration (write-up entry (vii) in Section 8) records the reason as cost. Partial hugo data exist and are labelled as such; they are excluded from all contrasts.
\item \emph{Run.} The remaining 11 instances were run (8 cli, 3 caddy). All 8 cli instances are pilot instances also used in Phase 1.
\end{itemize}

\subsection{Arms}\label{sec:arms}

The arms are defined as follows. Construction metadata for every arm and instance is in \texttt{harness/abl-*/contexts.json} and \texttt{harness/ablr-*/contexts.json}.

\begin{itemize}
\item \textbf{W}: the write set, that is, the code the gold patch edits, shown with full bodies.
\item \textbf{A1}: W plus the compiler-computed interfaces of every direct dependency of W.
\item \textbf{A1$-d$}: A1 with the interface of dependency group $d$ removed (leave-one-out).
\item \textbf{A1g}: W plus only the dependencies referenced by the gold patch. References are either a selector \texttt{alias.Name} resolved through W's imports, or a selector \texttt{v.Name} where \texttt{Name} is exported by exactly one direct dependency. This arm is an oracle by construction.
\item \textbf{A1t}: W plus only the dependencies referenced by the task text, under the task-text rule in force (Section~\ref{sec:taskrule}). When a predictor keeps no dependency, A1g or A1t is byte-identical to W.
\item \textbf{Wd} (Phase 3): the declarations of W reachable, through the type-checked reference graph, from identifiers used in the fail-to-pass tests and in task-text code spans or selectors. Each file keeps its package clause and imports, and omitted declarations are marked. When every declaration is reachable, Wd is byte-identical to W.
\item \textbf{A1r} (Phase 3): full W plus only the dependency declarations that W references, rendered without bodies.
\item \textbf{Wdr} (Phase 3): the slice Wd plus the dependency declarations that the slice references.
\item \textbf{Wpad} (Phase 3): W preceded by padding intended to match A1's prompt tokens. The padding is drawn from same-module packages outside W, its dependencies and dependents, in a fixed hash order. Package names are recorded in each \texttt{contexts.json} under \texttt{pad\_meta} (5 and 12 packages on the two short caddy instances). Matching held at the median (ratio 1.00) but fell short on caddy 30d3bd2a8959 (183,345 against 237,893 tokens, ratio 0.77) and 6727b2f3a92b (229,911 against 277,505, ratio 0.83); these two ratios are computed from the token counts in Table~\ref{tab:p3inst}. The pool was exhausted below A1's size on those instances (write-up entry (iv) in Section 8 of the preregistration). Wpad contrasts without these two instances were not computed. The ``inert'' property has not been verified.
\end{itemize}

Opus 5.5 also ran an arm labelled A0, which appears in Tables~\ref{tab:p2inst} and~\ref{tab:p2con}. The results file does not define A0, and no reason for running it only on Opus 5.5 is recorded. We exclude it from the analysis because its construction cannot be stated from the results file, so its contrasts cannot be interpreted. So that the exclusion is not selective, we report its results here. A0$-$W was $-0.140$ $[-0.440, +0.120]$, LB $-0.400$, signs 1/2/7, $n=10$. On Opus 5.5, A0 exceeded W on caddy 7910dd94eea2 (0.60 against 0.00, $k=5$) and equalled W on cli 2bcdfbb1ac02 (0.80). It fell below W on caddy 6727b2f3a92b and cli 20147bf81fed (0.00 against 1.00 on each).

\subsection{Models, samples and spend}

The models are Haiku 5.5 (\texttt{claude-haiku-5-5}), \texttt{claude-haiku-4-5-20251001} (labelled Haiku 4.5 in the tables), Sonnet 5.5 (\texttt{claude-sonnet-5-5}) and Opus 5.5 (\texttt{claude-opus-5-5}). Each sample is an independent single-shot call.

\paragraph{Phase 1.} Haiku 5.5 at $k=20$ on W, A1, A1g, A1t and every A1$-d$.

\paragraph{Phase 2.} The dense models ran W, A1, A1g and A1t:
\begin{itemize}
\item Haiku 5.5 at $k=100$;
\item \texttt{claude-haiku-4-5-20251001} at $k=20$;
\item Sonnet 5.5 at $k=10$.
\end{itemize}
Opus 5.5 ran W, A1t and the excluded A0 arm at $k=5$. Leave-one-out at $k=100$ on Haiku 5.5 formed stage 6 of Phase 2.

\paragraph{Phase 3.}
\begin{itemize}
\item Haiku 5.5 ran $k=50$ on W, A1, A1t, Wd, A1r, Wdr and Wpad. The W and A1 cells continue the $k=20$ screen samples, so each has 30 post-screen samples per instance.
\item Sonnet 5.5 ran $k=10$ on W, A1, Wd, A1r, Wdr and Wpad.
\end{itemize}

\paragraph{Spend.} Spend at list prices, in USD, was:
\begin{itemize}
\item Phase 1: USD 12.87 for 5020 samples;
\item Phase 2 screen: USD 1.58 for 162 samples;
\item Phase 2: USD 95.7 for 7754 samples;
\item Phase 3: USD 38.6 for 4751 samples;
\item total: USD 148.75 for 17687 samples.
\end{itemize}

By model, the spend was:
\begin{itemize}
\item \texttt{claude-haiku-5-5}: 16202 samples, USD 101.81;
\item \texttt{claude-haiku-4-5-20251001}: 520 samples, USD 7.32;
\item \texttt{claude-sonnet-5-5}: 860 samples, USD 31.28;
\item \texttt{claude-opus-5-5}: 105 samples, USD 8.34.
\end{itemize}

The by-model sample counts and costs both sum to the per-phase totals (17687 samples, USD 148.75).

For comparison, the pilot cost about \$73 at list prices. This figure comes from the erratum of 2026-10-12; the pilot paper states \$61.

\subsection{The task-text rule}\label{sec:taskrule}

\paragraph{Original rule (Phase 1).} Phase 1 used the rule in the original H5 paragraph. A dependency counted as task-referenced if any of the following appeared in the task text:
\begin{itemize}
\item a selector \texttt{alias.Name} resolved through W's imports, or \texttt{v.Name} where \texttt{Name} is exported by exactly one direct dependency;
\item a whole-word package name of length at least 4;
\item a unique exported identifier of length at least 6.
\end{itemize}

The Phase 1 reading recorded on 2026-10-12 found that the package-name clause matched English words (for example \texttt{context}, \texttt{source}, \texttt{text}, \texttt{status}, \texttt{extension}) rather than code references.

\paragraph{Amended rule (2026-10-12, fixed before any Phase 2 call).} The task-text predictor now counts a dependency only if the task text contains one of:
\begin{itemize}
\item a selector \texttt{alias.Name} resolving to it;
\item an exported identifier of that dependency inside a Markdown code span;
\item its import path.
\end{itemize}
Whole-word package names in prose no longer count. Phases 2 and 3 use the amended rule.

\subsection{Statistics}

Paired contrasts use the instance-level percentile bootstrap with 2,000 resamples and seed 2. For each contrast we report:
\begin{itemize}
\item the mean per-instance difference;
\item its two-sided 95\% interval;
\item the one-sided 95\% lower bound (LB);
\item paired sign counts ($+/-/0$);
\item the number of instances $n$.
\end{itemize}

The decision rules are as follows.
\begin{itemize}
\item \emph{Non-inferiority.} A contrast is called non-inferior at $\eps = 0.05$ only when LB $> -0.05$.
\item \emph{Unresolved differences.} An interval that contains zero is reported as an unresolved difference, not as equivalence.
\item \emph{Non-degeneracy.} Following the note of 2026-10-13, every Phase 2 and Phase 3 contrast is accompanied by the non-degeneracy condition (best arm $\geq 0.1$), and per-instance absolute rates are given in Tables~\ref{tab:p2inst} and~\ref{tab:p3inst}. The reason is that a relative contrast between arms that all fail carries no information about sufficiency.
\end{itemize}

Section 6 of the preregistration names both a two-sided 95\% bootstrap interval and one-sided $\alpha=0.05$. The analysis amendment of 2026-10-12 adopted the one-sided 95\% lower bound. Under the two-sided interval's lower end, the only verdict that changes is Haiku 5.5 Wd$-$W on the post-screen set ($-0.053$), which would not be non-inferior. The NI column in Table~\ref{tab:p3con} uses the one-sided bound.

For leave-one-out, the null reference draws every cell of an instance from one pooled rate. In Phase 1 it counts removals lowering and raising the rate by more than $\eps$ separately ($k=20$); in Phase 2 stage 6 it counts removals with $|$difference$| > 0.05$ ($k=100$). Both use 2,000 simulations.

Section 6 of the preregistration names a different seed for the paired bootstrap. The analysis amendment of 2026-10-12 records that paired differences use seed 2.

\subsection{Amendments and deviations as dated in the preregistration}

\begin{itemize}
\item \textbf{2026-10-12, pilot erratum.} One pilot instance had an empty task text and is invalid. The ablation study excludes it, and a non-empty task-text check was added to mining.
\item \textbf{2026-10-12, pilot cost erratum.} Pilot spend is about \$73, not \$61.
\item \textbf{2026-10-12, harness fault in Phase 1.} A build-environment fault invalidated verdicts for 4,456 samples. Of these, 4,409 were re-verified from stored responses with a patch-hash equality check, and 44 were re-sampled. The faulty aggregate was seen before the repair.
\item \textbf{2026-10-12, analysis amendment.} The amendment introduced the one-sided LB, sign counts, the estimate without the most influential instance, ceiling and floor counts and a null reference. It also recorded seed 2. It adopted the one-sided 95\% lower bound, the more lenient of the two readings of Section 6, after the exploratory Phase 1 run and after the faulty Phase 1 aggregate had been seen; the amendment describes this as no change to the decision rule, which is contestable because Section 6's H2 text names the two-sided interval. The only verdict in this report that depends on the choice is Haiku 5.5 Wd$-$W on the post-screen set.
\item \textbf{2026-10-12, Phase 1 reading.} This reading was recorded before any further call.
\item \textbf{2026-10-12, amendment for the next phase.} The amendment comprises:
  \begin{enumerate}
  \item[(i)] the task-text rule above;
  \item[(ii)] the W screen for Phase 2;
  \item[(iii)] estimating per-dependency necessity only at $k \geq 100$, with the null reference reported alongside;
  \item[(iv)] set-level precision and recall of the task-text predictor against the gold-referenced set, computed without model calls.
  \end{enumerate}
\item \textbf{2026-10-13, Phase 2 note.} This note was recorded after stages 2 to 5 and before any further call. It covers the floor problem, the non-degeneracy condition and the restriction of leave-one-out to three responsive instances. It also records the absence of bd81afda9ab6 and the replacement of docker/cli (no root \texttt{go.mod}) by restic and golangci-lint.
\item \textbf{2026-10-13, amendment for Phase 3.} This amendment was fixed before any Phase 3 call. It defines arms Wd, A1r and Wdr (H6, descriptive) and Wpad (H7). It also sets the $[0.1, 0.9]$ screen, relaxes the mining limits with re-examination of previously rejected commits, fixes the per-model arm sets and $k$, and requires same-day pairing within instance.
\item \textbf{Write-up entry, deviations recorded after the phases.} (i) amendment (iv) not performed; (ii) 8 of 11 Phase 3 instances are pilot instances; (iii) Phase 3 pairs were not same-day; (iv) Wpad pool exhausted on two caddy instances; (v) post-screen set designated primary after seeing data; (vi) one-sided bound adopted; (vii) H5b computable and reported, hugo dropped with partial data labelled; (viii) by-model spend rounding error corrected. Below, items of this entry are cited as ``write-up entry (i)'' and so on, to distinguish them from the items of the 2026-10-12 amendment.
\end{itemize}

Events absent from every entry of Section 8 are listed in Section~\ref{sec:lim}.

\section{Results}

\begin{table}[t]\centering\footnotesize
\caption{Phase 2: per-instance pass rates on the 10 screened instances (Haiku 5.5 $k{=}100$, Haiku 4.5 $k{=}20$, Sonnet 5.5 $k{=}10$, Opus 5.5 $k{=}5$). An arm byte-identical to another arm carries that arm's rate. ``deps'' is the number of direct dependencies; ``kept'' the dependencies retained by the gold-referenced (A1g) and task-referenced (A1t) predictors.}\label{tab:p2inst}
\begin{tabular}{llrrr p{3.1cm} p{2.9cm} p{2.9cm} p{2.6cm}}\toprule
instance & repo & deps & A1g & A1t & Haiku 5.5 & Haiku 4.5 & Sonnet 5.5 & Opus 5.5\\\midrule
4f45ded75556 & caddyserver & 8 & 1 & 0 & W\,0.03 A1\,0.03 A1g\,0.01 A1t\,0.03 & W\,0.00 A1\,0.00 A1g\,0.00 A1t\,0.00 & W\,0.00 A1\,0.00 A1g\,0.00 A1t\,0.00 & A0\,0.00 W\,0.00 A1t\,0.00\\
6727b2f3a92b & caddyserver & 13 & 0 & 0 & W\,0.44 A1\,0.22 A1g\,0.44 A1t\,0.44 & W\,0.00 A1g\,0.00 A1t\,0.00 & W\,0.20 A1\,0.10 A1g\,0.20 A1t\,0.20 & A0\,0.00 W\,1.00 A1g\,1.00 A1t\,1.00\\
7910dd94eea2 & caddyserver & 12 & 1 & 0 & W\,0.01 A1\,0.07 A1g\,0.02 A1t\,0.01 & -- & W\,0.00 A1\,0.00 A1g\,0.00 A1t\,0.00 & A0\,0.60 W\,0.00 A1t\,0.00\\
c8591a1e46e4 & caddyserver & 5 & 1 & 2 & W\,0.57 A1\,0.46 A1g\,0.49 A1t\,0.45 & W\,0.40 A1\,0.15 A1g\,0.15 A1t\,0.05 & W\,0.20 A1\,0.60 A1g\,0.00 A1t\,0.00 & A0\,0.00 W\,0.00 A1t\,0.00\\
015e6cf6ac39 & cli & 0 & 0 & 0 & W\,0.00 A1\,0.00 A1g\,0.00 A1t\,0.00 & W\,0.00 A1\,0.00 A1g\,0.00 A1t\,0.00 & W\,0.00 A1\,0.00 A1g\,0.00 A1t\,0.00 & A0\,0.00 W\,0.00 A1\,0.00 A1g\,0.00 A1t\,0.00\\
20147bf81fed & cli & 50 & 3 & 0 & W\,0.00 A1\,0.01 A1g\,0.00 A1t\,0.00 & W\,0.00 A1\,0.00 A1g\,0.00 A1t\,0.00 & W\,0.00 A1\,0.00 A1g\,0.00 A1t\,0.00 & A0\,0.00 W\,1.00 A1t\,1.00\\
2bcdfbb1ac02 & cli & 4 & 2 & 2 & W\,0.00 A1\,0.27 A1g\,0.14 A1t\,0.14 & W\,0.00 A1\,0.00 A1g\,0.00 A1t\,0.00 & W\,0.00 A1\,0.10 A1g\,0.10 A1t\,0.10 & A0\,0.80 W\,0.80\\
9b6585be5218 & cli & 8 & 1 & 0 & W\,0.01 A1\,0.00 A1g\,0.00 A1t\,0.01 & W\,0.00 A1\,0.00 A1g\,0.00 A1t\,0.00 & W\,0.10 A1\,0.00 A1g\,0.00 A1t\,0.10 & A0\,0.00 W\,0.00 A1t\,0.00\\
9ee0afb43414 & cli & 10 & 0 & 0 & W\,0.00 A1\,0.00 A1g\,0.00 A1t\,0.00 & W\,0.00 A1\,0.00 A1g\,0.00 A1t\,0.00 & W\,0.00 A1\,0.00 A1g\,0.00 A1t\,0.00 & A0\,0.00 W\,0.00 A1g\,0.00 A1t\,0.00\\
ff28e72b59be & cli & 62 & 0 & 0 & W\,0.00 A1\,0.00 A1g\,0.00 A1t\,0.00 & W\,0.00 A1\,0.05 A1g\,0.00 A1t\,0.00 & W\,1.00 A1\,0.80 A1g\,1.00 A1t\,1.00 & A0\,1.00 W\,1.00 A1g\,1.00 A1t\,1.00\\
\bottomrule\end{tabular}\end{table}

\begin{table}[t]\centering\small
\caption{Phase 2: paired contrasts on the 10 screened instances. Instance-level percentile bootstrap (2{,}000 resamples): mean difference with two-sided 95\% interval, one-sided 95\% lower bound (LB), sign counts $+/-/0$, and $n$. Degenerate instances (best arm $<0.1$): Haiku 5.5 7/10, Haiku 4.5 9/10, Sonnet 5.5 5/10, Opus 5.5 5/10.}\label{tab:p2con}
\begin{tabular}{llrrrr}\toprule
model & contrast & mean [95\% CI] & LB & $+/-/0$ & $n$\\\midrule
haiku-5-5 & A1t $-$ A1 & $+0.002$ $[-0.043, +0.059]$ & $-0.038$ & 2/4/4 & 10\\
haiku-5-5 & A1g $-$ A1 & $+0.004$ $[-0.042, +0.061]$ & $-0.035$ & 2/4/4 & 10\\
haiku-5-5 & W $-$ A1 & $+0.000$ $[-0.076, +0.071]$ & $-0.060$ & 3/3/4 & 10\\
haiku-5-5 & A1t $-$ W & $+0.002$ $[-0.036, +0.042]$ & $-0.024$ & 1/1/8 & 10\\
haiku-4-5-20251001 & A1t $-$ A1 & $-0.019$ $[-0.044, +0.000]$ & $-0.037$ & 0/2/6 & 8\\
haiku-4-5-20251001 & A1g $-$ A1 & $-0.006$ $[-0.019, +0.000]$ & $-0.019$ & 0/1/7 & 8\\
haiku-4-5-20251001 & W $-$ A1 & $+0.025$ $[-0.019, +0.094]$ & $-0.013$ & 1/1/6 & 8\\
haiku-4-5-20251001 & A1t $-$ W & $-0.039$ $[-0.117, +0.000]$ & $-0.117$ & 0/1/8 & 9\\
sonnet-5-5 & A1t $-$ A1 & $-0.020$ $[-0.170, +0.080]$ & $-0.150$ & 3/1/6 & 10\\
sonnet-5-5 & A1g $-$ A1 & $-0.030$ $[-0.180, +0.060]$ & $-0.150$ & 2/1/7 & 10\\
sonnet-5-5 & W $-$ A1 & $-0.010$ $[-0.120, +0.070]$ & $-0.100$ & 3/2/5 & 10\\
sonnet-5-5 & A1t $-$ W & $-0.010$ $[-0.060, +0.030]$ & $-0.050$ & 1/1/8 & 10\\
opus-5-5 & A1t $-$ W & $+0.000$ $[+0.000, +0.000]$ & $+0.000$ & 0/0/9 & 9\\
opus-5-5 & A0 $-$ W & $-0.140$ $[-0.440, +0.120]$ & $-0.400$ & 1/2/7 & 10\\
\bottomrule\end{tabular}\end{table}

\begin{table}[t]\centering\small
\caption{Leave-one-dependency-out at $k{=}100$ (Haiku 5.5) on the three instances with A1 $\geq 0.1$. Difference is pass(A1 minus group) $-$ pass(A1). The null reference draws every cell of an instance from one pooled rate and counts removals with $|\text{difference}| > 0.05$.}\label{tab:loo}
\begin{tabular}{llrrr}\toprule
instance (A1, W; tokens A1) & removed group & pass & difference & $n$\\\midrule
6727b2f3a92b (A1 0.22, W 0.44; 277,505) & caddyevents & 0.18 & $-0.04$ & 100\\
 & httpcaddyfile & 0.20 & $-0.02$ & 100\\
 & caddyconfig & 0.21 & $-0.01$ & 100\\
 & network & 0.27 & $+0.05$ & 100\\
 & headers & 0.28 & $+0.06$ & 100\\
 & v2 & 0.29 & $+0.07$ & 100\\
 & caddyfile & 0.30 & $+0.08$ & 100\\
 & cmd & 0.31 & $+0.09$ & 100\\
 & caddyhttp & 0.33 & $+0.11$ & 100\\
 & fileserver & 0.33 & $+0.11$ & 100\\
 & rewrite & 0.35 & $+0.13$ & 100\\
 & internal & 0.36 & $+0.14$ & 100\\
 & caddytls & 0.44 & $+0.22$ & 100\\
\multicolumn{5}{l}{\quad removals with $|$difference$|>0.05$: 9 of 13; null reference 5.5 [1, 12]; gold-referenced kept []; task-referenced kept []}\\\addlinespace
c8591a1e46e4 (A1 0.46, W 0.57; 84,817) & caddyfile & 0.34 & $-0.12$ & 100\\
 & caddyconfig & 0.42 & $-0.04$ & 100\\
 & httpcaddyfile & 0.45 & $-0.01$ & 100\\
 & v2 & 0.51 & $+0.05$ & 100\\
 & caddyhttp & 0.52 & $+0.06$ & 100\\
\multicolumn{5}{l}{\quad removals with $|$difference$|>0.05$: 2 of 5; null reference 2.3 [0, 5]; gold-referenced kept ['caddyhttp']; task-referenced kept ['caddyfile', 'caddyhttp']}\\\addlinespace
2bcdfbb1ac02 (A1 0.27, W 0.00; 16,293) & queries & 0.00 & $-0.27$ & 100\\
 & iostreams & 0.18 & $-0.09$ & 100\\
 & cmdutil & 0.22 & $-0.05$ & 100\\
 & client & 0.28 & $+0.01$ & 100\\
\multicolumn{5}{l}{\quad removals with $|$difference$|>0.05$: 2 of 4; null reference 1.5 [0, 4]; gold-referenced kept ['queries', 'iostreams']; task-referenced kept ['queries', 'iostreams']}\\\addlinespace
\bottomrule\end{tabular}\end{table}

\begin{table}[t]\centering\scriptsize\setlength{\tabcolsep}{3pt}
\caption{Phase 3: prompt tokens and pass rates on the 11 instances (Haiku 5.5 $k{=}50$; Sonnet 5.5 $k{=}10$). A trailing $=$ marks an arm byte-identical to another arm (it carries that arm's rate). ``slice'' is kept/total declarations of W; ``dep decls'' is referenced dependency declarations for the slice/for all of W.}\label{tab:p3inst}
\begin{tabular}{llrrrrrrllrrrrrrrrrrrrr}\toprule
instance & repo & tok W & tok Wd & tok A1 & tok A1r & tok Wdr & tok Wpad & slice & dep decls & H W & H Wd & H A1 & H A1r & H Wdr & H Wpad & H A1t & S W & S Wd & S A1 & S A1r & S Wdr & S Wpad\\\midrule
30d3bd2a8959 & caddyserver & 171,794 & 82,533 & 237,893 & 198,792 & 89,753 & 183,345 & 189/606 & 42/126 & 0.98 & 0.96 & 0.94 & 1.00 & 1.00 & 1.00 & 0.98= & 1.00 & 1.00 & 1.00 & 1.00 & 1.00 & 1.00\\
6727b2f3a92b & caddyserver & 158,565 & 73,923 & 277,505 & 203,263 & 93,699 & 229,911 & 181/459 & 68/157 & 0.34 & 0.20 & 0.26 & 0.06 & 0.18 & 0.06 & 0.34= & 0.30 & 0.20 & 0.20 & 0.30 & 0.40 & 0.00\\
c8591a1e46e4 & caddyserver & 12,137 & 10,047 & 84,817 & 17,460 & 13,474 & 84,798 & 23/37 & 23/27 & 0.64 & 0.38 & 0.36 & 0.52 & 0.30 & 0.36 & 0.44 & 0.20 & 0.00 & 0.50 & 0.00 & 0.00 & 0.30\\
2bcdfbb1ac02 & cli & 2,367 & 2,367 & 16,293 & 4,443 & 4,443 & 16,341 & 7/7 & 14/14 & 0.00 & 0.00= & 0.22 & 0.60 & 0.60= & 0.00 & --= & 0.30 & 0.30= & 0.10 & 0.00 & 0.00= & 0.20\\
326faaac8b5c & cli & 21,011 & 21,011 & 67,579 & 31,357 & 31,357 & 68,049 & 44/44 & 79/79 & 0.86 & 0.86= & 0.84 & 0.94 & 0.94= & 0.96 & 0.86= & 1.00 & 1.00= & 1.00 & 1.00 & 1.00= & 1.00\\
3f6a16a9f8c7 & cli & 13,485 & 4,045 & 24,999 & 17,612 & 5,744 & 24,960 & 6/17 & 10/43 & 0.88 & 0.94 & 0.84 & 0.88 & 0.74 & 0.74 & 0.88= & 1.00 & 1.00 & 1.00 & 1.00 & 1.00 & 1.00\\
8f7afbdbdbd6 & cli & 16,439 & 16,439 & 56,667 & 21,820 & 21,820 & 56,531 & 23/23 & 44/44 & 0.98 & 0.98= & 0.92 & 0.62 & 0.62= & 0.28 & 0.46 & 1.00 & 1.00= & 1.00 & 1.00 & 1.00= & 1.00\\
ad2a33837712 & cli & 78,155 & 16,657 & 83,776 & 79,540 & 16,776 & 83,757 & 110/401 & 2/19 & 0.46 & 0.60 & 0.58 & 0.46 & 0.74 & 0.40 & 0.46= & 1.00 & 1.00 & 1.00 & 0.70 & 0.60 & 1.00\\
b87e58a1d390 & cli & 26,184 & 6,235 & 68,609 & 32,648 & 6,434 & 68,718 & 14/61 & 3/70 & 0.42 & 0.46 & 0.48 & 0.58 & 0.42 & 0.50 & 0.42= & 1.00 & 0.90 & 0.60 & 0.70 & 0.60 & 0.80\\
babd691446b9 & cli & 22,957 & 6,347 & 53,549 & 24,996 & 7,779 & 53,884 & 19/99 & 13/25 & 0.34 & 0.60 & 1.00 & 1.00 & 1.00 & 0.16 & --= & 1.00 & 1.00 & 1.00 & 1.00 & 1.00 & 1.00\\
cba9646b8c25 & cli & 86,253 & 32,808 & 103,793 & 90,653 & 36,818 & 102,489 & 134/420 & 30/38 & 0.80 & 0.86 & 0.80 & 0.80 & 0.72 & 0.82 & --= & 1.00 & 0.90 & 0.90 & 0.80 & 1.00 & 1.00\\
\bottomrule\end{tabular}\end{table}

\begin{table}[t]\centering\small
\caption{Phase 3: paired contrasts on the 11 instances. Haiku 5.5 is reported with W and A1 restricted to post-screen samples ($n{=}30$ per instance; other arms $n{=}50$) and with all samples; Sonnet 5.5 at $k{=}10$. NI: one-sided 95\% lower bound $> -0.05$. Ratio: median per-instance prompt-token ratio.}\label{tab:p3con}
\begin{tabular}{l rrrr r rrr rrr}\toprule
 & \multicolumn{5}{c}{Haiku 5.5, post-screen} & \multicolumn{3}{c}{Haiku 5.5, all samples} & \multicolumn{3}{c}{Sonnet 5.5}\\\cmidrule(lr){2-6}\cmidrule(lr){7-9}\cmidrule(lr){10-12}
contrast & mean [95\% CI] & LB & NI & $+/-/0$ & ratio & mean [95\% CI] & LB & $+/-/0$ & mean [95\% CI] & LB & $+/-/0$\\\midrule
Wd $-$ W & $+0.009$ $[-0.053, +0.072]$ & $-0.043$ & yes & 4/4/3 & 0.47 & $+0.013$ $[-0.064, +0.089]$ & $-0.051$ & 5/3/3 & $-0.045$ $[-0.082, -0.009]$ & $-0.082$ & 0/4/7\\
A1r $-$ A1 & $+0.021$ $[-0.095, +0.141]$ & $-0.079$ & no & 5/4/2 & 0.48 & $+0.020$ $[-0.089, +0.129]$ & $-0.071$ & 6/3/2 & $-0.073$ $[-0.182, +0.009]$ & $-0.164$ & 2/4/5\\
Wdr $-$ A1 & $+0.002$ $[-0.102, +0.116]$ & $-0.082$ & no & 4/5/2 & 0.27 & $+0.002$ $[-0.091, +0.105]$ & $-0.078$ & 4/6/1 & $-0.064$ $[-0.191, +0.036]$ & $-0.164$ & 2/3/6\\
Wdr $-$ W & $+0.042$ $[-0.133, +0.234]$ & $-0.107$ & no & 6/5/0 & 0.52 & $+0.051$ $[-0.129, +0.251]$ & $-0.102$ & 5/5/1 & $-0.109$ $[-0.218, -0.018]$ & $-0.200$ & 1/4/6\\
A1 $-$ W & $+0.039$ $[-0.076, +0.185]$ & $-0.061$ & no & 3/4/4 & 2.33 & $+0.049$ $[-0.067, +0.202]$ & $-0.053$ & 4/6/1 & $-0.045$ $[-0.145, +0.045]$ & $-0.127$ & 1/4/6\\
Wpad $-$ W & $-0.138$ $[-0.281, -0.018]$ & $-0.255$ & no & 3/7/1 & 2.35 & $-0.129$ $[-0.273, -0.011]$ & $-0.245$ & 4/6/1 & $-0.045$ $[-0.118, +0.009]$ & $-0.100$ & 1/3/7\\
Wpad $-$ A1 & $-0.178$ $[-0.373, -0.024]$ & $-0.335$ & no & 5/6/0 & 1.00 & $-0.178$ $[-0.371, -0.020]$ & $-0.336$ & 4/6/1 & $+0.000$ $[-0.064, +0.064]$ & $-0.055$ & 3/2/6\\
A1r $-$ W & $+0.060$ $[-0.114, +0.248]$ & $-0.086$ & no & 5/6/0 & 1.28 & $+0.069$ $[-0.100, +0.260]$ & $-0.080$ & 5/3/3 & $-0.118$ $[-0.200, -0.045]$ & $-0.191$ & 0/5/6\\
A1t $-$ A1 & $-0.046$ $[-0.177, +0.051]$ & $-0.156$ & no & 2/2/4 & 0.66 & $-0.047$ $[-0.180, +0.045]$ & $-0.155$ & 5/3/0 & $+0.083$ $[+0.000, +0.217]$ & $+0.000$ & 2/0/4\\
A1t $-$ W & $-0.087$ $[-0.223, +0.000]$ & $-0.203$ & no & 0/2/6 & 1.00 & $-0.090$ $[-0.220, +0.000]$ & $-0.195$ & 0/2/6 & $+0.000$ $[+0.000, +0.000]$ & $+0.000$ & 0/0/6\\
\bottomrule\end{tabular}\end{table}


\subsection{Phase 1: leave-one-dependency-out on the pilot instances}

All Phase 1 figures are reproduced from the frozen \texttt{harness/results-abl/ABLATION-REPORT.md}. Phase 1 has no table in this report.

\paragraph{Shift relative to the pilot.} A1 on Haiku 5.5 at $k=20$ differed from A1 in the pilot at $k=10$ by $-0.041$ $[-0.074, -0.013]$ over 23 instances. This shift (which also confounds $k$ and date) is close in size to $\eps$.

\paragraph{Predictor arms against A1.}
\begin{itemize}
\item \textbf{A1g$-$A1}: $+0.017$ $[-0.002, +0.041]$, LB $+0.000$, signs 5/2/16. A1g is therefore non-inferior to A1 at $\eps$. Without the most influential instance (2bcdfbb1ac02) the estimate is $+0.011$ $[-0.005, +0.032]$. Median tokens were 13,513 against 39,586 for A1.
\item \textbf{A1t$-$A1}: $+0.007$ $[-0.020, +0.030]$, LB $-0.015$, signs 5/3/15. A1t is non-inferior at $\eps$. Without 2bcdfbb1ac02 the estimate is $+0.000$ $[-0.023, +0.023]$. Median tokens were 15,335.
\end{itemize}
Both intervals contain zero. These are non-inferiority readings, not evidence of equal performance.

\paragraph{W against A1.} W$-$A1 was $-0.028$ $[-0.096, +0.017]$, LB $-0.083$, signs 4/4/15, so W is not non-inferior to A1. Without babd691446b9 the estimate becomes $+0.000$ $[-0.025, +0.023]$. That is a post-hoc single-instance exclusion, which we do not treat as a finding.

\paragraph{How many instances could separate the arms.} The predictors kept no dependency on 14/23 instances (A1g) and 11/23 (A1t); on those instances the arm equals W. On 8 instances W and A1 were both at ceiling, and on 5 both were at floor. W was inexpressible on 10. Most paired differences were exactly zero (16 for A1g$-$A1 and 15 for A1t$-$A1), so only a minority of instances carry the contrasts.

\paragraph{Necessity counts.} On the 17/23 informative instances (A1 pass $> 0$) there were 135 dependency groups.
\begin{itemize}
\item Removal lowered the pass rate by more than $\eps$ for 33/135 groups and raised it by more than $\eps$ for 33/135.
\item The null reference gives 30.8 $[13, 49]$ lowering and 30.4 $[12, 53]$ raising. Both observed counts lie inside these ranges.
\item The bootstrap fraction 0.244 $[0.178, 0.319]$ ignores clustering by instance and should not be read as an estimate of necessity.
\end{itemize}

Per instance (necessary/helps/groups), the counts are concentrated. Examples are 326faaac8b5c 0/14/16, b87e58a1d390 11/2/15, 8f7afbdbdbd6 9/2/11 and 2bcdfbb1ac02 0/3/4. With $k=20$ per cell, however, these patterns are compatible with sampling noise.

Predictor precision and recall against this label are printed in the results file. As recorded on 2026-10-12, they are not interpretable as predictor quality because the label itself is uninformative. The slack ratio tokens(A1)/tokens(W $\cup$ necessary) has mean 2.71 $[1.66, 4.06]$ and median 1.56, and it inherits the same problem.

\subsection{Phase 2: four models on ten screened instances}

Table~\ref{tab:p2inst} gives the per-instance rates and Table~\ref{tab:p2con} the paired contrasts. A dash in Table~\ref{tab:p2inst} marks cells with no reported rate.

\paragraph{Prompt sizes.} Median prompt tokens on the 10 instances were 21,946 for W, 84,817 for A1, and 49,822 for each of A1g and A1t.

\paragraph{Degenerate instances.} The best arm was below 0.1 on:
\begin{itemize}
\item 7/10 instances for Haiku 5.5;
\item 9/10 for \texttt{claude-haiku-4-5-20251001};
\item 5/10 for Sonnet 5.5;
\item 5/10 for Opus 5.5.
\end{itemize}
The relative contrasts therefore rest on few informative instances and are close to vacuous.

\paragraph{Haiku 5.5 ($k=100$, $n=10$).}
\begin{itemize}
\item A1t$-$A1: $+0.002$ $[-0.043, +0.059]$, LB $-0.038$, signs 2/4/4.
\item A1g$-$A1: $+0.004$ $[-0.042, +0.061]$, LB $-0.035$, signs 2/4/4.
\item W$-$A1: $+0.000$ $[-0.076, +0.071]$, LB $-0.060$.
\item A1t$-$W: $+0.002$ $[-0.036, +0.042]$, LB $-0.024$, signs 1/1/8.
\end{itemize}
The lower bounds of A1t$-$A1, A1g$-$A1 and A1t$-$W clear $-0.05$. A1t equals W on 8/10 instances, so A1t$-$W is zero by construction there. Only three instances meet the non-degeneracy condition, so we do not read these as evidence of sufficiency.

\paragraph{\texttt{claude-haiku-4-5-20251001} ($k=20$).} This model has no cells on caddy 7910dd94eea2 and no A1 cell on caddy 6727b2f3a92b. Counting byte-identical arms once, the shown cells account for 22 cells, or 440 samples; the 80 excess samples equal the four missing cells (three distinct arms on 7910dd94eea2 and A1 on 6727b2f3a92b), which indicates these cells were sampled and not reported. The reason for their exclusion is not recorded in the results file. Section 6 of the preregistration requires missing data to be recorded and the instance flagged; neither is done in the results file, so this is listed as a deviation in Section~\ref{sec:lim}.

Its contrasts against A1 therefore use $n=8$, and A1t$-$W uses $n=9$. All contrasts have intervals touching or containing zero:
\begin{itemize}
\item A1t$-$A1: $-0.019$ $[-0.044, +0.000]$;
\item A1g$-$A1: $-0.006$ $[-0.019, +0.000]$;
\item W$-$A1: $+0.025$ $[-0.019, +0.094]$;
\item A1t$-$W: $-0.039$ $[-0.117, +0.000]$, LB $-0.117$.
\end{itemize}
With 9/10 instances degenerate, these numbers carry essentially no information.

\paragraph{Sonnet 5.5 ($k=10$, $n=10$).} No contrast was resolved:
\begin{itemize}
\item A1t$-$A1: $-0.020$ $[-0.170, +0.080]$;
\item A1g$-$A1: $-0.030$ $[-0.180, +0.060]$;
\item W$-$A1: $-0.010$ $[-0.120, +0.070]$;
\item A1t$-$W: $-0.010$ $[-0.060, +0.030]$, LB $-0.050$ (at the margin as reported to three decimals; the rule requires LB $> -0.05$). As for Haiku 5.5, A1t$-$W is zero by construction on the 8 instances where A1t equals W.
\end{itemize}

\paragraph{Opus 5.5 ($k=5$).} Opus had no A1 arm; the A1 entry on 015e6cf6ac39 (0 dependencies) is W by identity. Opus also has no A1t cell on cli 2bcdfbb1ac02, although A1t there keeps 2 of the 4 dependencies and so differs from W; the cause is not recorded. A1t$-$W was $+0.000$ $[+0.000, +0.000]$ with signs 0/0/9. This is a degenerate interval, not evidence of sufficiency.

\paragraph{Cells with large differences between arms.} The averages hide cells with large differences between arms, in both directions:
\begin{itemize}
\item On cli 2bcdfbb1ac02, Haiku 5.5 passed 0/100 with W and 27/100 with A1. A1g and A1t, which are identical here, passed 0.14.
\item On caddy c8591a1e46e4, Sonnet 5.5 passed 6/10 with A1 against 2/10 with W, and 0/10 with each of A1g and A1t. Haiku 5.5 on the same instance passed 0.57 with W and 0.46 with A1 at $k=100$, which is not a resolved difference. \texttt{claude-haiku-4-5-20251001} passed 0.40 with W and 0.15 with A1 at $k=20$.
\item On caddy 6727b2f3a92b, Haiku 5.5 passed 44/100 with W and 22/100 with A1. Here W had 158,565 prompt tokens and A1 had 277,505. In Phase 3 post-screen samples this cell did not replicate (W 0.27, A1 0.27).
\end{itemize}

\paragraph{Cross-model differences.} Some instances that were degenerate for one model were solved by another:
\begin{itemize}
\item On cli ff28e72b59be, Haiku 5.5 passed 0.00 in every arm. Sonnet 5.5 passed 1.00 with W, and Opus 5.5 passed 1.00 in every analysed arm.
\item On cli 20147bf81fed, Opus 5.5 passed 1.00 with W.
\end{itemize}

\paragraph{Unverified diagnostic.} A diagnostic was recorded on 2026-10-13 during a harness session. It is not a script output. It examined three cells: caddy c8591a1e46e4 on Sonnet 5.5, cli 2bcdfbb1ac02 on Haiku 5.5 and caddy 6727b2f3a92b on Haiku 5.5. In those cells, the passing patches introduced no code identifier that only the dependency documentation supplied. The only dependency-only tokens in passing patches were English words from comments. This observation is unverified by the analysis scripts.

\subsection{Leave-one-out at $k=100$}

Table~\ref{tab:loo} gives the results. Each count of removals with absolute difference above 0.05 lies within its simulated null range:
\begin{itemize}
\item caddy 6727b2f3a92b: 9 of 13 against 5.5 $[1, 12]$;
\item caddy c8591a1e46e4: 2 of 5 against 2.3 $[0, 5]$;
\item cli 2bcdfbb1ac02: 2 of 4 against 1.5 $[0, 4]$.
\end{itemize}
On caddy c8591a1e46e4 and cli 2bcdfbb1ac02 the null range spans every possible count, so the count test could not have rejected. No necessity claim is supported at the level of counts.

\paragraph{caddy 6727b2f3a92b.} Here A1 is 0.22 and W is 0.44. Most removals raised the pass rate; only removing caddyevents, httpcaddyfile or caddyconfig lowered it. Removing caddytls gave 0.44 ($+0.22$), equal to W. Neither predictor kept any dependency on this instance.

\paragraph{caddy c8591a1e46e4.} Here A1 is 0.46 and W is 0.57. Removing caddyfile gave 0.34 ($-0.12$), the largest drop. The gold-referenced predictor kept caddyhttp, and the task-referenced predictor kept caddyfile and caddyhttp.

\paragraph{cli 2bcdfbb1ac02.} Here A1 is 0.27 and W is 0.00. The individual removals gave:
\begin{itemize}
\item queries: 0.00 ($-0.27$);
\item iostreams: 0.18 ($-0.09$);
\item cmdutil: 0.22 ($-0.05$);
\item client: 0.28 ($+0.01$).
\end{itemize}
Both predictors kept exactly queries and iostreams. The corresponding Phase 2 cell (A1g $=$ A1t) passed 0.14. We report the single cell for queries descriptively: with 100 samples per arm, a rate of 0.27 falling to 0.00 is unlikely to be sampling noise alone. Even so, it is one cell on one instance.

\subsection{Phase 3: declaration-level arms and matched padding}

Table~\ref{tab:p3inst} gives prompt tokens and per-instance rates. The Haiku W and A1 columns pool the $k=20$ screen and 30 post-screen samples ($n=50$); the primary contrasts in Table~\ref{tab:p3con} use the post-screen 30. Table~\ref{tab:p3con} gives the paired contrasts.

\paragraph{Primary sample set.} We take as primary the Haiku 5.5 set in which W and A1 are restricted to the post-screen samples ($n=30$ per instance; other arms $n=50$). This choice was not preregistered. We made it because the screen samples were used to select the instances: including them risks regression to the mean in the selected cells. The per-instance screen-to-post-screen changes in Section~\ref{sec:lim} show such movement. Both sets are reported.

\paragraph{Non-degeneracy and identical arms.} Every instance met the non-degeneracy condition for both models (Haiku 11/11, Sonnet 11/11). Sonnet 5.5 had both W and A1 at 1.00 on 6/11 instances, which limits every Sonnet contrast.

Wd was byte-identical to W on cli 2bcdfbb1ac02, 326faaac8b5c and 8f7afbdbdbd6, where every declaration was reachable. The Haiku A1t arm was byte-identical to another arm on most instances. On babd691446b9, 2bcdfbb1ac02 and cba9646b8c25, Table~\ref{tab:p3inst} marks A1t as identical to another arm (\texttt{--=}) without a rate or the identity of the duplicated arm; the results file has the same entries. On 2bcdfbb1ac02 the marking conflicts with Phase 2, where A1t kept queries and iostreams, 2 of the 4 dependencies, and so differed from both W and A1.

\paragraph{A1t on Haiku 5.5.} Haiku 5.5 has A1t rates on 8 of the 11 instances. On 6 of these A1t is byte-identical to another arm, and its rate equals W's on each of them, so on those 6 A1t$-$W is zero by construction and A1t$-$A1 equals W$-$A1. Distinct A1t samples exist only on c8591a1e46e4 (0.44) and 8f7afbdbdbd6 (0.46). On the post-screen set, A1t$-$A1 was $-0.046$ $[-0.177, +0.051]$, LB $-0.156$, signs 2/2/4, $n=8$: not non-inferior. A1t$-$W was $-0.087$ $[-0.223, +0.000]$, LB $-0.203$, signs 0/2/6, carried entirely by those two instances. On all samples the figures are $-0.047$ $[-0.180, +0.045]$, LB $-0.155$, signs 5/3/0; and $-0.090$ $[-0.220, +0.000]$, LB $-0.195$.

Sonnet 5.5 did not run A1t: its Phase 3 arm list under the 2026-10-13 amendment excludes it. The Sonnet 5.5 entries in the two A1t rows of Table~\ref{tab:p3con} therefore do not come from A1t samples. They are W$-$A1 and W$-$W imputed through byte identity on 6 instances, and should be read as absent. In particular, the Sonnet A1t$-$A1 row is not a test of H5b.

\paragraph{Prompt sizes.} Median prompt tokens were:
\begin{itemize}
\item W 22,957;
\item Wd 16,439;
\item A1 68,609;
\item A1r 31,357;
\item Wdr 16,776;
\item Wpad 68,718.
\end{itemize}

\paragraph{Haiku 5.5, primary (post-screen) set.}
\begin{itemize}
\item \textbf{Wd$-$W}: $+0.009$ $[-0.053, +0.072]$, one-sided LB $-0.043$, signs 4/4/3, median token ratio 0.47. This meets the margin at $\eps = 0.05$ under the one-sided bound adopted in the analysis amendment, but not under the lower end of the two-sided interval ($-0.053$); the interval contains zero. On 3 of the 11 instances Wd is byte-identical to W, so 3 of the paired differences are zero by construction. On the 8 instances where Wd differs from W the signs are 4/4/0; the mean and interval on those 8 instances were not computed.
\item \textbf{A1r$-$A1}: $+0.021$ $[-0.095, +0.141]$, LB $-0.079$, ratio 0.48. Not non-inferior.
\item \textbf{Wdr$-$A1}: $+0.002$ $[-0.102, +0.116]$, LB $-0.082$, ratio 0.27. Not non-inferior.
\item \textbf{Wdr$-$W}: $+0.042$ $[-0.133, +0.234]$, LB $-0.107$. Not non-inferior.
\item \textbf{A1$-$W}: $+0.039$ $[-0.076, +0.185]$, ratio 2.33. Unresolved.
\item \textbf{A1r$-$W}: $+0.060$ $[-0.114, +0.248]$. Unresolved.
\item \textbf{Wpad$-$W}: $-0.138$ $[-0.281, -0.018]$, LB $-0.255$, signs 3/7/1, ratio 2.35. The interval excludes zero.
\item \textbf{Wpad$-$A1}: $-0.178$ $[-0.373, -0.024]$, LB $-0.335$, signs 5/6/0, median ratio 1.00; matching failed on caddy 30d3bd2a8959 (0.77) and 6727b2f3a92b (0.83). The interval excludes zero.
\end{itemize}

\paragraph{Haiku 5.5, all samples.}
\begin{itemize}
\item Wd$-$W was $+0.013$ $[-0.064, +0.089]$ with LB $-0.051$, so it is not non-inferior on this set. The same 3 differences are zero by construction (signs 5/3/3). The non-inferiority verdict for Wd$-$W therefore depends on the sample set.
\item The padding contrasts were essentially unchanged: Wpad$-$W $-0.129$ $[-0.273, -0.011]$ and Wpad$-$A1 $-0.178$ $[-0.371, -0.020]$.
\item A1$-$W was $+0.049$ $[-0.067, +0.202]$ with LB $-0.053$.
\end{itemize}

\paragraph{Sonnet 5.5 ($k=10$; W $=$ A1 $=$ 1.00 on 6/11).}
\begin{itemize}
\item \textbf{Wd$-$W}: $-0.045$ $[-0.082, -0.009]$, LB $-0.082$, signs 0/4/7. Not non-inferior. The point estimate is negative and the instance-level interval excludes zero (see the caveat below).
\item \textbf{A1r$-$W}: $-0.118$ $[-0.200, -0.045]$, signs 0/5/6. The instance-level interval excludes zero.
\item \textbf{Wdr$-$W}: $-0.109$ $[-0.218, -0.018]$. The instance-level interval excludes zero.
\item \textbf{A1r$-$A1}: $-0.073$ $[-0.182, +0.009]$. Unresolved.
\item \textbf{Wdr$-$A1}: $-0.064$ $[-0.191, +0.036]$. Unresolved.
\item \textbf{A1$-$W}: $-0.045$ $[-0.145, +0.045]$. Unresolved.
\item \textbf{Wpad$-$W}: $-0.045$ $[-0.118, +0.009]$. Unresolved, and not evidence of a padding effect.
\item \textbf{Wpad$-$A1}: $+0.000$ $[-0.064, +0.064]$, LB $-0.055$. Unresolved and not non-inferior.
\end{itemize}
On Sonnet 5.5 ($k=10$), the point estimates of Wd$-$W, A1r$-$W and Wdr$-$W were negative and their instance-level intervals exclude zero. Each rests on 4--5 instances with differences of 1--4 samples out of 10. The interval does not include within-cell sampling variance, so these are indications, not resolved effects.

\paragraph{Instances that carry the Haiku contrasts.} The instance-level rates in Table~\ref{tab:p3inst} show which instances carry the Haiku contrasts.
\begin{itemize}
\item \emph{Padding.} Wpad fell well below W on:
  \begin{itemize}
  \item babd691446b9 (0.16 against 0.34, with A1 1.00);
  \item 8f7afbdbdbd6 (0.28 against 0.98);
  \item 6727b2f3a92b (0.06 against 0.34).
  \end{itemize}
  On 2bcdfbb1ac02, Wpad passed 0.00 against 0.22 for A1. On c8591a1e46e4, Wpad and A1 both passed 0.36 at nearly equal prompt sizes (84,798 and 84,817 tokens).
\item \emph{Slice and dependency arms.} The slice and dependency-declaration arms gained on some instances and lost on others:
  \begin{itemize}
  \item babd691446b9: Wd 0.60 and A1r 1.00, against W 0.34;
  \item ad2a33837712: Wdr 0.74, against W 0.46;
  \item 8f7afbdbdbd6: A1r 0.62, against A1 0.92;
  \item c8591a1e46e4 and 6727b2f3a92b: the slice losses discussed in Section~\ref{sec:lim}.
  \end{itemize}
\end{itemize}

\section{Reading of the evidence}

This section separates what the data show from interpretation. Paragraphs marked \emph{Interpretation} go beyond the measurements.

\subsection{What is supported and what is not}

\paragraph{Supported, within the exploratory scope.}
\begin{enumerate}
\item An individual Haiku 5.5 cell at $k=100$ showed a large difference between arms: on cli 2bcdfbb1ac02, W passed 0/100 and A1 27/100. On the same instance, removing queries from A1 removed all passes.
\end{enumerate}

\paragraph{Fragile.}
\begin{enumerate}
\item On Haiku 5.5 in Phase 3, Wpad had lower mean pass rates than W ($-0.138$ $[-0.281,-0.018]$, signs 3/7/1) and A1 ($-0.178$ $[-0.373,-0.024]$, signs 5/6/0), driven by large drops on 2--4 instances; token matching fell short on two caddy instances, the ``inert'' property of the padding is unverified, the intervals ignore within-cell variance and are not adjusted for the ten contrasts, and the percentile bootstrap at $n=11$ may under-cover.
\item Wd met the margin against W only on the post-hoc post-screen set and only under the one-sided bound; it did not with screen samples included, nor under the two-sided bound, nor on Sonnet 5.5.
\item On Sonnet 5.5 ($k=10$), the point estimates of Wd$-$W, A1r$-$W and Wdr$-$W were negative and their instance-level intervals exclude zero. Each rests on 4--5 instances with differences of 1--4 samples out of 10. The interval does not include within-cell sampling variance, so these are indications, not resolved effects.
\item On caddy 6727b2f3a92b, Haiku 5.5 passed 44/100 with W and 22/100 with A1 in Phase 2. The difference did not replicate in Phase 3 post-screen samples (W 0.27, A1 0.27).
\end{enumerate}

\paragraph{Not supported.}
\begin{itemize}
\item Any claim that dependency interfaces are generally unnecessary once W is present: Phase 1 W$-$A1 is not non-inferior.
\item Any claim that particular dependencies are necessary at the count level: every leave-one-out count lies inside its null range, and on two of the three instances the count test could not have rejected.
\item Any claim of sufficiency on Phase 2 floor instances.
\item Any claim that the static predictors recover sufficient context.
\end{itemize}

A defensible summary of Phase 2 is the following. Across ten screened tasks from two repositories we detected no clear average advantage from adding dependency interfaces. Individual task--model comparisons showed both benefits and harms, and seven tasks were near floor on the primary model.

\subsection{Lower pass rates under padding on headroom instances, and the flat pilot curve}

\paragraph{Data.} On Haiku 5.5, Wpad$-$W was $-0.138$ at a median token ratio of 2.35. Wpad$-$A1 was $-0.178$ at a median ratio of 1.00, with token matching failed on two caddy instances. A1$-$W, which adds about as many tokens of dependency interfaces as Wpad adds padding at the median (ratio 2.33), was $+0.039$ (unresolved).

\paragraph{Interpretation.} Taken together, these contrasts suggest that for Haiku 5.5 the harm is not a function of prompt length alone. On these instances, at one length, padding was worse than dependency interfaces of similar size; whether this conflicts with reports that length alone hurts even with non-informative filler \cite{du2025} is not tested here.

Two confounds remain. First, if A1, like Wpad, places W after the added text, then Wpad$-$A1 holds W's position fixed. Wpad$-$W changes length and W's distance from the prompt start together. Position \cite{liu2024} is therefore separated from content only in Wpad$-$A1, and only if A1 serialises W last; the serialisation order of A1 was not checked for this report. Second, the Sonnet 5.5 contrasts are unresolved and close to ceiling, so we make no claim about Sonnet.

\paragraph{Interpretation (pilot curve).} The pilot's padding contrast did not resolve a difference. Phase 1 offers one plausible reason. On the pilot instances, W and A1 were both at ceiling on 8 instances and both at floor on 5, so most instances had no room to degrade. The Phase 3 screen admitted only instances with a rate in $[0.1, 0.9]$, and lower pass rates under padding were observed there for Haiku 5.5. This reading is post hoc and has not been tested with a load dose curve. Wpad is matched to A1 rather than to the preregistered load levels, so it is not a test of H1.

\subsection{Declaration-level slicing}

\paragraph{Data.} On Haiku 5.5, Wd's one-sided lower bound against W clears $-0.05$ on the primary set only (not under the lower end of the two-sided interval), at about half the tokens. A1r and Wdr do not meet the margin against A1, at ratios of 0.48 and 0.27. On Sonnet 5.5, Wd, Wdr and A1r each had lower point estimates than W (Wd$-$W $-0.045$, Wdr$-$W $-0.109$, A1r$-$W $-0.118$), subject to the caveat on within-cell sampling variance above.

\paragraph{Interpretation.} The slice is a candidate for reducing context within the 0.05 margin for one model on these instances. The evidence is fragile: the verdict changes with the sample set and with the choice of bound, and it does not hold for the second model. We do not read Wd as $\eps$-sufficient in general.

\subsection{Granularity versus package selection, and possible non-additivity on cli 2bcdfbb1ac02}

\paragraph{Data.} On cli 2bcdfbb1ac02 with Haiku 5.5:
\begin{itemize}
\item W passed 0.00 and A1 passed 0.27 (Phase 2, $k=100$).
\item Removing queries alone gave 0.00, and removing iostreams alone gave 0.18.
\item Removing cmdutil alone gave 0.22, and removing client alone gave 0.28.
\item Yet the arm that keeps exactly queries and iostreams (A1g $=$ A1t) passed 0.14, about half of A1's rate.
\item In Phase 3, A1r (W plus only the dependency declarations W references, without bodies; 4,443 tokens against 16,293 for A1) passed 0.60 against 0.22 for A1.
\end{itemize}

\paragraph{Interpretation.} The data suggest that removal effects may not be additive on this instance. Under additivity, removing both cmdutil ($-0.05$) and client ($+0.01$) would be expected to give a rate only slightly below A1's 0.27; the observed rate of the arm that keeps only queries and iostreams was 0.14. The gap between 0.14 and the rate expected under additivity is itself not resolved at $k=100$. The cells also come from different runs, and a drift of this size has been observed on this instance (A1 0.30 in the Phase 3 screen and 0.17 post-screen), so this is not established. Leave-one-out measures the effect of removing one group while all others remain, and cannot identify a minimal sufficient subset when groups interact or are redundant.

A1r (Phase 3) passed 0.60, against 0.14 for A1g $=$ A1t in Phase 2. The comparison crosses phases, and A1r is smaller than A1. This is one instance, so granularity and size are not separated. On Sonnet 5.5 the same A1r passed 0.00 against 0.30 for W.

\subsection{The cross-model asymmetry}

\paragraph{Data.} The same context had different effects on different models:
\begin{itemize}
\item On caddy c8591a1e46e4 in Phase 2, the point estimates favoured A1 on Sonnet 5.5 (6/10 against 2/10) and W on both Haiku models (Haiku 5.5 0.57 against 0.46 at $k=100$; Haiku 4.5 0.40 against 0.15 at $k=20$). None is resolved at these $k$.
\item In Phase 3 the signs of the slice and dependency contrasts differ between Haiku 5.5 and Sonnet 5.5. Padding lowered Haiku 5.5 pass rates with intervals excluding zero; the Sonnet 5.5 padding contrasts were unresolved.
\item Two instances at floor for Haiku 5.5 were solved by another model (cli ff28e72b59be by Sonnet 5.5 W and Opus 5.5; cli 20147bf81fed by Opus 5.5 W).
\end{itemize}

Sonnet's W cell on cli 2bcdfbb1ac02 was 0.00 in Phase 2 and 0.30 in Phase 3, both at $k=10$, which is compatible with sampling noise. For Haiku 5.5, A1 differed by $-0.041$ between the pilot and Phase 1, a shift that also confounds $k$ and date.

\paragraph{Interpretation.} Whether a context is sufficient, and whether added text helps or harms, appears to depend on the reader. Several of the differences above are unresolved, so this remains a hypothesis. This is the sense of ``reader effects'' in the title. The diagnostic of 2026-10-13 found no dependency-only code identifiers in passing patches in the three cells it examined. That diagnostic is a harness-session observation, not a script output. The attribution recorded in the preregistration on the same date (load, framing, semantic reminders) is an untested mechanism, and we do not claim it is shown.

\subsection{The predictors}

\paragraph{Data.} On the pilot instances the predictors kept nothing on 14/23 (A1g) and, under the original task-text rule, 11/23 (A1t). Precision and recall against the Phase 1 necessity label are uninterpretable, as recorded on 2026-10-12. In Phase 2 the predictors behaved as follows:
\begin{itemize}
\item On caddy 6727b2f3a92b, both predictors kept nothing. This matched the Phase 2 Haiku 5.5 cell (W 0.44 against A1 0.22) but not the Phase 3 post-screen cell (0.27 each).
\item On caddy c8591a1e46e4, Sonnet 5.5 passed 0.60 with A1 and 0.00 with either predictor arm.
\item On cli 2bcdfbb1ac02, the predictors kept the two groups whose individual removal hurt most, and the arm still passed only 0.14.
\end{itemize}

\paragraph{Interpretation.} The tested static selectors missed at least one clear dependency-sensitive case, so their ability to recover sufficient context remains unestablished. On floor instances, a predictor that keeps nothing cannot be credited with a correct selection, because every arm fails. The set-level comparison of task-text and gold-referenced sets in amendment (iv) was not computed (see Section~\ref{sec:lim}).

\subsection{What this does and does not say about static computability}

The Stage 1 protocol asks whether a context computed by the compiler, without model calls, is $\eps$-sufficient. These phases do not answer that question. They show that one statically computed context (Wd) met the margin against W for one model on one sample set and under one choice of bound, and failed it for another model. The package-level predictors (A1g, A1t) did not reproduce full A1's pass rate on the two Phase 2 cells with the largest A1-over-W point estimates (cli 2bcdfbb1ac02 at $k=100$: 0.14 against 0.27; caddy c8591a1e46e4 on Sonnet 5.5 at $k=10$, unresolved: 0.00 against 0.60). The declaration-level A1r (Phase 3) exceeded A1 on cli 2bcdfbb1ac02 for Haiku 5.5 but scored 0.00 against 0.30 for W on Sonnet 5.5.

\emph{Interpretation.} If the unverified diagnostic is right, the differences between arms need not involve identifiers that only dependency text supplies; whether they reflect presentation rather than information is untested. A purely structural criterion of sufficiency would then need to be paired with a model-specific check. This is a hypothesis for the confirmatory study, not a result.

\section{Limitations and deviations}\label{sec:lim}

\paragraph{Data-integrity events.} The following events are recorded, together with their consequences.

\begin{itemize}
\item \textbf{2026-10-12, Phase 1.}
  \begin{itemize}
  \item A build-environment fault invalidated verdicts for 4,456 samples. Of these, 4,409 were re-verified from stored responses with a patch-hash equality check, and 44 were re-sampled.
  \item The re-verified and re-sampled counts do not sum to the invalidated count, and the results file does not record the disposition of the difference.
  \item The faulty aggregate was seen before the repair. Phase 1 remains exploratory for this reason among others.
  \end{itemize}
\item \textbf{2026-10-12/13, Phase 2.}
  \begin{itemize}
  \item A disk-full interruption contaminated the hugo baseline and 17 mining rejections, which were purged and re-examined.
  \item docker/cli was replaced by restic and golangci-lint because it has no root \texttt{go.mod}.
  \item Instance cli bd81afda9ab6 is absent because its arms failed to build.
  \end{itemize}
\item \textbf{2026-10-13, Phase 2 stage 6.} Leave-one-out was restricted to three responsive instances. Completed samples on floor instances were kept. The restriction followed the note of 2026-10-13, which limited leave-one-out to instances with A1 $\geq 0.1$ on Haiku 5.5. It forfeits the ability to detect rescue effects, in which removing a distracting dependency raises a floor instance.
\item \textbf{2026-10-14, Phase 3.}
  \begin{itemize}
  \item The shared Go build cache reached about 150 GB, caused by per-patch relinking of test binaries, and lanes were restarted with per-lane caches. The results file separately attributes a 182 GB build cache to the hugo instance; which cache that figure refers to is not stated in the results file, so the two figures are not reconciled here.
  \item The hugo instance was dropped.
  \item Seventeen records written while the disk was full were purged together with their cached-verdict dependants and re-sampled. The originals are retained by the author.
  \end{itemize}
\end{itemize}
The two counts of 17 refer to different kinds of record (mining rejections on 2026-10-12/13, sample records on 2026-10-14) and are listed as separate events in the results file.

\paragraph{Deviations recorded in the write-up entry of Section 8.} The following deviations are recorded in the write-up entry listed in the amendments above:
\begin{itemize}
\item \textbf{Wpad token matching.} Wpad was not token-matched to A1 on caddy 30d3bd2a8959 and 6727b2f3a92b, contrary to the 2026-10-13 amendment, because the same-module padding pool was exhausted below A1's size on those instances (write-up entry (iv)).
\item \textbf{Same-day pairing.} Pairs were not same-day: W and A1 began with the 2026-10-13 screen and the other arms ran on 2026-10-13/14 (write-up entry (iii)).
\item \textbf{The meaning of ``the Phase 3 instances'' in H5.} The H5 paragraph (2026-10-12) calls H5 confirmatory on ``the Phase 3 instances''. That phrase predates the Phase 3 defined on 2026-10-13. Because 8 of the 11 Phase 3 instances are pilot instances (write-up entry (ii)), on which H5 is exploratory by the preregistration, the phrase can only coherently refer to the confirmatory study. This report adopts that reading, under which H5 remains to be tested in the confirmatory study. The write-up entry (vii) records H5b as computable on the Phase 3 instances and reported there. Under that other reading, H5a was untested, since there was no A1g arm, and H5b would be a confirmatory test that failed (A1t$-$A1 LB $-0.156$ on $n=8$, six of them by byte identity).
\item \textbf{Amendment (iv).} Amendment (iv) (2026-10-12), set-level precision and recall of the task-text predictor on every mineable instance, was not computed (write-up entry (i)).
\item \textbf{Phase 3 sample set.} The choice of the post-screen sample set as primary for Phase 3 was made after the data were seen and was not preregistered (write-up entry (v)).
\end{itemize}

\paragraph{Events not in any entry of Section 8.} The results file states that all events are recorded in Section 8 of \texttt{preregistration.md}, but the version of that section used to prepare this report does not contain the following items. The archive is therefore inconsistent. The items are:
\begin{itemize}
\item \textbf{2026-10-12/13.} The hugo baseline contamination and the purge and re-examination of 17 mining rejections.
\item \textbf{2026-10-14.} The build-cache incident and the purge and re-sampling of 17 records.
\item \textbf{2026-10-14, relaxed mining.} The 2026-10-13 amendment relaxed the mining limits (at most 4 packages, at most 300 changed lines) and required re-examination of commits rejected only by the old limits. On 2026-10-14 the relaxed limits produced no new valid instance on restic and golangci-lint. The relaxed mining of caddy, hugo and cli was stopped by the disk incident and not rerun.
\item \textbf{Missing Phase 2 cells.} \texttt{claude-haiku-4-5-20251001} has no reported cells on caddy 7910dd94eea2 and no reported A1 cell on caddy 6727b2f3a92b, and Opus 5.5 has no A1t cell on cli 2bcdfbb1ac02. The sample count for \texttt{claude-haiku-4-5-20251001} indicates that the four missing cells of that model were sampled and not reported. No cause is recorded.
\item \textbf{Phase 2 screen sample count.} The 162 screen samples recorded are fewer than 35 instances at $k=5$ would give, and the shortfall is not explained.
\end{itemize}

\paragraph{Selection and $k$.} All Phase 2 and Phase 3 instances were selected by screens on the same model, and the screened population is an enriched subset of the mining population. All 8 cli instances of Phase 3 are pilot instances that had already been run on Haiku 5.5 in Phase 1, which adds to the selection effect.

The Phase 2 screen at $k=5$ is coarse: an instance with a modest true rate can return zero passes often. That screen also admitted floor instances, 7/10 on Haiku 5.5.

The Phase 3 screen used $k=20$, and the post-screen rates moved, in some cases substantially:
\begin{itemize}
\item 6727b2f3a92b: W 0.45 to 0.27;
\item c8591a1e46e4: A1 0.45 to 0.30;
\item 2bcdfbb1ac02: A1 0.30 to 0.17;
\item 3f6a16a9f8c7: W 0.75 to 0.97.
\end{itemize}

The Phase 1 necessity estimates at $k=20$ carry no information beyond the null reference. Leave-one-out at $k=100$ covered only three instances.

\paragraph{The hugo instance.} The one hugo instance selected in Phase 3 was dropped as unrunnable on the available disk (the preregistration records the reason as cost). The hugo repository is therefore unrepresented in Phases 2 and 3, apart from screening.

\paragraph{Two repositories.} All run instances come from \texttt{cli/cli} and \texttt{caddyserver/caddy}: 10 in Phase 2 and 11 in Phase 3, with some overlap. Repository-specific code structure may drive every effect reported. Repeated samples on the same instance do not create independent tasks.

\paragraph{Sonnet 5.5 at $k=10$.} Sonnet rates move in steps of 0.1, so one sample changes a cell visibly. With 6/11 Phase 3 instances at ceiling for W and A1, most Sonnet contrasts rest on a handful of instances and a few samples per cell. The Phase 2 Sonnet contrasts share the same granularity.

\paragraph{Slice losses on caddy c8591a1e46e4 and 6727b2f3a92b.}
\begin{itemize}
\item \emph{c8591a1e46e4.} The slice kept 23 of 37 declarations (10,047 against 12,137 tokens). Haiku 5.5 passed 0.38 with Wd against 0.64 with W, and Sonnet 5.5 passed 0.00 against 0.20.
\item \emph{6727b2f3a92b.} The slice kept 181 of 459 declarations. Haiku 5.5 passed 0.20 against 0.34, and Sonnet 5.5 passed 0.20 against 0.30.
\end{itemize}
The reachability rule, the marking of omitted declarations or the omission itself may be responsible; these were not separated.

\paragraph{Other limitations.} Several further limitations apply:
\begin{itemize}
\item A1g is an oracle, because it uses the gold patch.
\item Opus 5.5 had no A1 arm and ran at $k=5$.
\item The shift on Haiku 5.5 between the pilot and Phase 1 ($-0.041$), which also confounds $k$ and date, is close to $\eps$, which limits the meaning of a margin of 0.05 across days. Phase 3 pairs were not same-day: W and A1 began with the 2026-10-13 screen and the other arms ran on 2026-10-13/14 (write-up entry (iii)).
\item Token matching of Wpad failed on two caddy instances because the same-module pool was exhausted (write-up entry (iv)), and the ``inert'' property of the padding is unverified.
\end{itemize}

\section{Consequences for the registered report and future work}

\paragraph{Headroom screening for H1.} Lower pass rates under padding were observed for Haiku 5.5 on instances selected for headroom; the pilot padding contrast on unselected instances was unresolved. The confirmatory H1 instance set should therefore be screened so that both the unpadded arm and its padded versions can move. Selection should rely on separate screen samples, and the evaluation should use fresh samples, with the screened and mining populations reported separately. A triage screen followed by further screening of borderline instances would reduce misclassification.

\paragraph{Instances used here.} The pilot instances, including all 8 cli instances of Phase 3, and the caddy and other cli instances screened and run in these phases have been seen on the evaluated models; they should be excluded from the confirmatory set.

\paragraph{Declaration-level arms as candidate confirmatory arms.} Wd and A1r are candidates for confirmatory arms, with the token ratios reported alongside every non-inferiority reading. Given the negative Sonnet 5.5 point estimates and the dependence of the Haiku 5.5 verdict on the sample set and the choice of bound, they should be tested as hypotheses, not adopted as defaults.

\paragraph{A load dose curve.} Wpad tests a single load matched to A1. A dose curve over several padding levels, with W in a fixed position and with padding placed both before and after W, would separate length, content and position. This is what H1 already specifies for A1.

\paragraph{A second dense model.} The cross-model asymmetry means that conclusions from Haiku 5.5 alone may not transfer. A second dense model at a $k$ comparable to Haiku's would allow per-instance contrasts with useful resolution.

\paragraph{Instance scarcity.} Instance directories number 14 for caddy and 30 for cli, the latter including the pilot and earlier template/test instances, with single instances elsewhere. Relaxed mining of caddy, hugo and cli has not been rerun. The confirmatory sample size in the preregistration therefore depends on adding repositories and completing the relaxed mining. Build-cache isolation per lane should be in place from the start.

\paragraph{The obligation-graph idea.} An idea raised in a critique received by the author (not archived) is to make explicit, as a graph, the declarations a patch must be consistent with, and to derive context from that graph. It is untested here. The critical test is whether such a graph can be built before seeing the solution; a graph derived from the gold patch is an oracle diagnostic only. An explicit dependency would not imply that the dependency must be included. Structural relevance, reader need and presentation would need to be measured separately.

\section{Availability}
This report, its generated results file, tables and the verdict records are archived at Zenodo under DOI \url{https://doi.org/10.5281/zenodo.23460566}.


All data except the invalid originals, which are retained by the author, together with the harness code and analysis scripts, are in the repository \url{https://github.com/Shockware-Limited/suffbench-registered-report}. The materials are organised as follows:
\begin{itemize}
\item verdict records: \texttt{harness/results-abl}, \texttt{harness/results-p2}, \texttt{harness/results-p2screen} and \texttt{harness/results-p3};
\item arm metadata: \texttt{harness/abl-*/contexts.json} and \texttt{harness/ablr-*/contexts.json};
\item rejected commits with reasons: \texttt{harness/instances/rejected.jsonl}.
\end{itemize}

\texttt{analysis/phase23.py} generates the Phase 2 and Phase 3 numbers; Phase 1 numbers are from \texttt{analysis/ablation.py} via the frozen report \texttt{harness/results-abl/ABLATION-REPORT.md}; the 2026-10-13 diagnostic is not script output. \texttt{analysis/phase23\_tables.py} generates every table. The pilot erratum figures are taken from the preregistration.

\section*{Acknowledgements and AI-use disclosure}

The author used Anthropic's Claude Fable 5.1, which the author refers to as Lemma, through GitHub Copilot to design and run the phases and to write the analysis scripts. The author used Claude Opus 5.5 through the Anthropic API to draft this report and to review it adversarially against the generated results file. An adversarial critique of the Phase 1 result was obtained from a Claude agent run by the author (not archived). An external critique of the Phase 2 summary was also received (not archived). Both critiques are reflected in the text. The author is responsible for all claims.

\begin{thebibliography}{9}

\bibitem{definitions}
J.~C. Johansen.
Definitions note on $\eps$-sufficient context.
Shockware Limited, Zenodo.

\bibitem{pilot}
J.~C. Johansen.
Pilot study of $\eps$-sufficient context for language-model code editing in Go.
Shockware Limited, Zenodo.
DOI \url{https://doi.org/10.5281/zenodo.23307204}.
Erratum of 2026-10-12.

\bibitem{stage1}
J.~C. Johansen.
Interfaces or bodies? A test of $\eps$-sufficient context for LM code editing in Go.
Registered report, Stage 1 protocol. Shockware Limited.

\bibitem{joren2025}
Joren et al.
Sufficient Context: A New Lens on Retrieval Augmented Generation Systems.
In \emph{International Conference on Learning Representations (ICLR)}, 2025.

\bibitem{du2025}
Du et al.
Context Length Alone Hurts LLM Performance Despite Perfect Retrieval.
In \emph{Findings of the Association for Computational Linguistics: EMNLP}, 2025.

\bibitem{liu2024}
Liu et al.
Lost in the Middle: How Language Models Use Long Contexts.
\emph{Transactions of the Association for Computational Linguistics}, 2024.

\end{thebibliography}

\end{document}\begin{thebibliography}{9}

\bibitem{definitions}
J.~C. Johansen.
$\eps$-Sufficient Context for Bounded Readers: Definitions and a Non-Monotonicity Result for Code Editing.
Version 1.1, Shockware Limited, Zenodo, October 2026.
DOI \url{https://doi.org/10.5281/zenodo.23282824} (concept DOI 10.5281/zenodo.23282195).

\bibitem{pilot}
J.~C. Johansen.
Interfaces or Bodies? A 24-Instance Pilot Measurement of $\eps$-Sufficient Context for Language-Model Code Editing in Go.
Pilot report, version 1.0, Shockware Limited, Zenodo, October 2026.
DOI \url{https://doi.org/10.5281/zenodo.23307204}. Errata of 2026-10-12 are recorded in the preregistration.

\bibitem{stage1}
J.~C. Johansen.
Interfaces or Bodies? A Preregistered Test of $\eps$-Sufficient Context for Language-Model Code Editing in Go.
Registered report, Stage 1 protocol, Shockware Limited, Zenodo, October 2026.
DOI \url{https://doi.org/10.5281/zenodo.23307547}.

\bibitem{joren2025}
H.~Joren, J.~Zhang, C.-S. Ferng, D.-C. Juan, A.~Taly, and C.~Rashtchian.
Sufficient Context: A New Lens on Retrieval Augmented Generation Systems.
In \emph{International Conference on Learning Representations (ICLR)}, 2025.

\bibitem{du2025}
Y.~Du, M.~Tian, S.~Ronanki, S.~Rongali, S.~B. Bodapati, A.~Galstyan, A.~Wells, R.~Schwartz, E.~A. Huerta, and H.~Peng.
Context Length Alone Hurts LLM Performance Despite Perfect Retrieval.
In \emph{Findings of the Association for Computational Linguistics: EMNLP 2025}, 2025.

\bibitem{liu2024}
N.~F. Liu, K.~Lin, J.~Hewitt, A.~Paranjape, M.~Bevilacqua, F.~Petroni, and P.~Liang.
Lost in the Middle: How Language Models Use Long Contexts.
\emph{Transactions of the Association for Computational Linguistics}, 12:157--173, 2024.

\bibitem{gloaguen2026}
T.~Gloaguen, N.~M\"undler-Sasahara, M.~N. M\"uller, V.~Raychev, and M.~Vechev.
Evaluating AGENTS.md: Are Repository-Level Context Files Helpful for Coding Agents?
arXiv:2602.11988, 2026.

\end{thebibliography}

\end{document}
